\begin{trapbox}
\textbf{\textcolor{CTrap}{特殊三角形的坐标与边长不能混用}}

写 $I=(r,r)$ 的前提是：直角顶点在原点，两直角边沿正坐标轴。若平移、旋转图形，
应相应改变内心坐标，不能照抄 $(r,r)$。等腰直角三角形中，直角边 $t$ 与斜边 $w=\sqrt2t$
对应不同的半径公式；等边三角形三个角都是 $60^\circ$，不存在“直角等边三角形”。

\medskip
\textbf{\textcolor{CTrap}{切线段长的两个提醒}}

$x=AE=AF$ 是从 $A$ \textbf{沿边到切点}的长度，不是半径 $r$，也不是 $AI$。
要用 $x=p-a$，其中 $a=BC$ 是 $A$ 的对边；不能把它误写成 $p-AB$。
同理，$y$ 对应 $B$ 的对边 $CA$，$z$ 对应 $C$ 的对边 $AB$。

\medskip
\textbf{\textcolor{CTrap}{易错点一：到三边等距，不等于到三个顶点等距}}

\PIncenterFigure{2}{4}{6}{1}{a}{b}{c}{1}{蓝色实线垂直于各底边，是高；灰色虚线连到顶点，不是这些高。}
内心 $I$ 到三边的\textbf{垂直距离}相等，即 $ID=IE=IF$；
一般没有 $IA=IB=IC$。到三个顶点距离相等的是\textbf{外心}，不能混用。

\medskip
\textbf{\textcolor{CTrap}{易错点二：漏掉面积公式中的 $\frac12$}}

三个小三角形面积分别是 $ar/2,br/2,cr/2$，所以
\[
 S=\frac12r(a+b+c)=rp,\qquad r=\frac{2S}{a+b+c}.
\]
当使用\textbf{半周长 $p$} 时，$r=S/p$；当使用\textbf{整个周长 $L$} 时，$r=2S/L$。
两种写法只是分母表示不同。

\medskip
\textbf{\textcolor{CTrap}{易错点三：把连接内心到顶点的线段当作高}}

以 $BC$ 为底的 $\triangle IBC$，高是 $ID$，\textbf{不是} $IB$ 或 $IC$；
其他两块同理。只有垂线段才能当相应的高。

\medskip
\textbf{\textcolor{CTrap}{易错点四：只给“到三条边所在直线等距”就断定是内心}}

\begin{center}
\begin{tikzpicture}[scale=.75,font=\small,line join=round]
\coordinate (A) at (0,3); \coordinate (B) at (0,0); \coordinate (C) at (4,0);
\coordinate (I) at (1,1); \coordinate (J) at (6,6);
\draw[gray,dashed] (0,3)--(0,6.7) (4,0)--(6.7,0);
\draw[thick] (A)--(B)--(C)--cycle;
\draw[CStatement,thick] (I) circle[radius=1];
\draw[CStatement] (I)--(1,0) (I)--(0,1) (I)--(.4,1.8);
\draw[CTrap,thick,dashed] (J)--node[midway,right] {$6$} (6,0)
 (J)--node[midway,above] {$6$} (0,6)
 (J)--node[pos=.55,right=4pt] {$6$} (2.4,1.2);
\draw (5.75,0)--(5.75,.25)--(6,.25);
\draw (0,5.75)--(.25,5.75)--(.25,6);
\coordinate (H) at (2.4,1.2);
\pic[draw,angle radius=3mm] {right angle=J--H--C};
\foreach \p/\where in {A/left,B/below left,C/below,I/above right,J/above right}
 {\fill (\p) circle[radius=1.5pt];\node[\where=3pt] at (\p) {$\p$};}
\node[anchor=west,align=left] at (8,5.2) {$I$ 在内部\\[4pt]到三边距离均为 $1$};
\node[anchor=west,align=left,text=CTrap] at (8,2.2) {$J$ 在外部\\[4pt]到三边所在直线\\[4pt]距离均为 $6$};
\node[align=center] at (6,-1) {同一个 $3$、$4$、$5$ 三角形：等距还要看点是否在内部};
\end{tikzpicture}
\end{center}

三角形\textbf{外部}也可能有到三条边所在直线距离相等的点；
只有再加“\textbf{点在三角形内部}”这一条件，才能据此判定它是内心。
作图题尤其要区分内角平分线与外角平分线。

\medskip
\textbf{\textcolor{CTrap}{易错点五：题目数据本身可能互相矛盾}}

原例题写“周长 $24$、内切圆半径 $3$”，直接套式子得到
$S=3\times12=36$，\textbf{但这是不存在的三角形，不能当成真实例题}。
下面给出一个只用高中知识的检查：
设半周长 $p=12$，由海伦公式，
\[
 S^2=12(12-a)(12-b)(12-c).
\]
令 $x=12-a,y=12-b,z=12-c$，则 $x,y,z>0$，且
\[
 x+y+z=36-(a+b+c)=12.
\]
三个正数的和固定为 $12$ 时，乘积在三数相等时最大（基本不等式），所以
\[
 xyz\le4^3=64,\qquad S^2\le12\times64=768,
 \qquad S\le16\sqrt3<36.
\]
于是与 $S=36$ 矛盾。例2现改为\textbf{$r=2,L=24,S=24$}，
并用 $6$、$8$、$10$ 直角三角形验证其确实存在。

\medskip
\textbf{\textcolor{CTrap}{易错点六：三边不能围成三角形}}

例如 $2,3,5$，虽然看起来给出了三边长度，
但 $2+3=5$，三点落在同一直线上，没有真正的内切圆。
使用本结论前必须确认任意两边之和\textbf{严格大于}第三边。

\medskip
\textbf{\textcolor{CTrap}{易错点七：看见三边就机械套海伦公式}}

若三边满足勾股定理，\textbf{先用直角三角形面积}通常更省步骤。
例如 $5,12,13$，直接 $S=5\times12/2=30$，
再代 $r=2S/(5+12+13)=2$。
\end{trapbox}
